\chapter{需求分析}
\section{使用场景与功能需求}
学生首先选择日期、座位与学习时段，系统返回预约结果及记录编号。学生需要提前离开时，可以取消本人预约；取消记录保留在数据库中，便于确认状态变化。查询功能读取仍处于有效状态的预约，并统计指定日期的预约数量与总时长。原型中的用户与座位由初始化程序建立。

主要功能及其约束见表\ref{tab:requirements}。
\begin{table}[H]
  \centering
  \caption{主要功能需求}
  \label{tab:requirements}
  \tablefont
  \begin{tabularx}{\textwidth}{@{}p{24mm}p{48mm}Y@{}}
    \toprule
    功能 & 输入或对象 & 处理要求\\
    \midrule
    创建预约 & 用户、座位、日期与时段 & 输入有效，座位与用户均无重叠预约\\
    取消预约 & 用户与预约编号 & 仅允许取消本人有效预约\\
    查询预约 & 日期 & 返回该日有效记录与累计预约时长\\
    状态保留 & 已取消记录 & 保留原记录，将其排除于占用判断\\
    \bottomrule
  \end{tabularx}
\end{table}

\section{业务规则}
\subsection{时段与冲突规则}
原型设置每日开放时间为 08:00 至 22:00，以当天零时起的分钟数存储起止时刻。结束时刻必须晚于开始时刻，预约整体必须位于开放范围内。同一日期中，一个座位在重叠时段只能对应一条有效预约；同一用户也不能同时预约多个座位。前一预约的结束时刻与后一预约的开始时刻相同，视为相邻时段，可分别预约。

\subsection{取消与权限规则}
取消操作同时核对记录编号、用户编号和当前状态。属于其他用户或已经取消的记录，均返回操作失败。成功取消后，记录状态变为已取消，其时段立即从冲突查询中排除。重复取消不会进一步修改数据。

\section{质量需求}
预约创建需要保证检查与写入的连续性，避免两个请求同时依据过期查询结果作出决定。数据层应检查用户、座位之间的引用关系，业务层应返回可区分的结果码。测试需要包含成功路径与边界输入，并提供清晰的运行结果，便于定位规则与实现之间的差异。
